% Generated draft: requires mathematical and source review.
\documentclass[11pt]{article}
\usepackage[margin=25mm]{geometry}
\usepackage{fontspec}
\setmainfont{Latin Modern Roman}
\usepackage{amsmath,amssymb,mathtools,unicode-math}
\setmathfont{Latin Modern Math}
\usepackage{enumitem,needspace,xurl,hyperref}
\setlength{\parindent}{0pt}
\setlength{\parskip}{5pt}
\setlength{\emergencystretch}{3em}
\newcommand{\N}{\mathbb{N}}
\newcommand{\Z}{\mathbb{Z}}
\newcommand{\Q}{\mathbb{Q}}
\newcommand{\R}{\mathbb{R}}
\newcommand{\C}{\mathbb{C}}
\newcommand{\F}{\mathbb{F}}
\newcommand{\A}{\mathbb{A}}
\newcommand{\PP}{\mathbb{P}}
\newcommand{\E}{\mathbb{E}}
\newcommand{\eps}{\varepsilon}
\newcommand{\dd}{\,\mathrm{d}}
\begin{document}
\section*{UnsolvedMath ID: 30001815}
\textbf{Source problem code:} catalogue entry
\subsection{Definitions and mathematical statement}
Let \(G\) be a semisimple Lie group, meaning that its Lie algebra is semisimple. A subgroup \(\Gamma\leq G\) is cocompact and discrete if it is discrete in the topology of \(G\) and the quotient \(G/\Gamma\) is compact. A strict total order \(<\) on a group \(H\) is left invariant if
\[
\forall a,b,c\in H,\qquad a<b\Longrightarrow ca<cb.
\]
The group \(H\) is left-orderable if it admits such an order. A group \(\Gamma\) is virtually left-orderable if some subgroup \(H\leq\Gamma\) with \([\Gamma:H]<\infty\) is left-orderable.

The supplied question is whether there exist a semisimple Lie group \(G\) and a cocompact arithmetic subgroup \(\Gamma\leq G\) such that
\[
\forall H\leq\Gamma,\qquad
[\Gamma:H]<\infty\Longrightarrow
H\text{ admits no left-invariant strict total order}.
\]
Here an arithmetic subgroup is understood to be discrete, with the arithmetic condition interpreted according to the cited report. The supplied record does not specify the algebraic-group data defining arithmeticity, the permitted compact factors, or conventions concerning connectedness and the centre of \(G\). These details require source review; no additional restrictions on \(G\) are imposed in this draft.

The formulation is provisional because the supplied curation describes the extracted passage as truncated immediately before an explicitly numbered Open problem 1, while separate verification metadata asserts that the displayed question was checked against the official report. These supplied claims do not resolve the provenance issue within the available evidence.
\subsection{Short English statement}
Does there exist a cocompact arithmetic subgroup of a semisimple Lie group that is not virtually left-orderable? Every subgroup of finite index must fail to admit a left-invariant total order. The supplied open label and verification claims require independent review, particularly because the curation also reports a truncated source passage.
\subsection{Sources}
\begin{itemize}
\item Supplied local UnsolvedMath registry record, permanent ID \(30001815\), source problem code \texttt{catalogue entry}, title \emph{Virtual Nonorderability of Cocompact Arithmetic Groups}, in the Oberwolfach Reports Open Problems collection. The record assigns proposed year 2011 and names no individual proposer.
\item \emph{Arithmetic Groups vs. Mapping Class Groups: Similarities, Analogies and Differences} (2011), Oberwolfach Reports, pp. 1637--1708, DOI \texttt{10.4171/owr/2011/30}. Supplied source citation and literature-source link labelled ``original Oberwolfach report'': \url{https://doi.org/10.4171/owr/2011/30}.
\item Literature assessment dated exactly \texttt{2026-08-21}: the supplied curation labels the problem open and says that the extracted passage is truncated immediately before an explicitly numbered Open problem 1. It reports insufficient claim text to establish a later resolution or disproof. This is supplied dated curation, not an independently verified literature assessment; the record requests independent verification of time-sensitive status and cited claims.
\item The supplied statement status is \texttt{corrected\_verified}. Its accompanying verification text asserts comparison with the official report and normalization of notation, context, and question form without changing the source-level claim. This assertion is reproduced as source metadata and has not been independently verified here.
\end{itemize}
\end{document}
